% A10: common represented-full-exact rigid groupoid, AP/CCO no SAP.
\section{Full exact functional correspondences and rigid groupoids}
\label{sec:foundational-rigid-groupoid}

Work in a homogeneous complete finite binary language with finitely
many named, individually transitive original sorts, symmetric
same-sort orbitals, and converse-paired mixed orbitals. Assume
semantic homomorphic CCO, but not SAP. Products are \emph{actual
represented} relational composition, including logical identity.
Exact means full identity-inclusive typed word-exact as in
\cref{thm:full-typed-exact-balls}.

\begin{definition}[Full exact bridge]
\label{def:foundational-exact-bridge}
Given a full exact invariant equivalence $K_a$ on sort $a$,
a \emph{bridge} to sort $b$ is a nonempty full exact relation
$R:a\to b$ satisfying
\[
 K_aR=R,\qquad RR^\smile=K_a.
 \tag{FB}\label{eq:exact-bridge}
\]
Its target kernel is $K_b=R^\smile R$.
The identity of the quotient object $(a,K_a)$ is $K_a$,
not necessarily original real equality.
\end{definition}

\begin{theorem}[Full exact rigid-bridge groupoid]
\label{thm:foundational-rigid-groupoid}
Starting with one $(a_*,K_*)$ and closing under bridges, converse
and composition gives a \emph{finite connected groupoid} of
full exact quotient-functional correspondences, with:
\begin{enumerate}[label=(\roman*),leftmargin=*]
\item the target kernel $K_b$ determined by $b$, independently
of which bridge path was used;
\item full exact arrows $R:a\to b$ with
$K_aR=R=RK_b$, $RR^\smile=K_a$,
$R^\smile R=K_b$, inducing actual equivariant bijections
$M_a/K_a\to M_b/K_b$;
\item elementary abelian $2$-groups as isotropy groups,
with each nonempty hom-set a torsor under either endpoint group;
\item freeness: two different arrows $a\to b$ cannot agree
at even one quotient point.
\end{enumerate}
No finite-index parity pivot or distinguished phase is assumed.
\end{theorem}
\begin{proof}
The relation $R^\smile R$ is symmetric, full exact by convolution,
and reflexive by surjectivity of nonempty invariant $R$ on a
transitive target sort. Since
\[
 (R^\smile R)^2=R^\smile K_aR=R^\smile R,\qquad
 R(R^\smile R)=K_aR=R,
\]
it is an exact equivalence and $R$ is saturated at both ends.
The two equations $RR^\smile=K_a$ and $R^\smile R=K_b$
make the induced relation on quotient classes the graph of a
total bijection; equivariance is automatic from invariance.

If $R,S:a_*\to b$ have a common source kernel, the same-sort
invariant relation $T=R^\smile S$ is symmetric. Therefore
\[
 R^\smile R=R^\smile SS^\smile R=TT^\smile
          =T^\smile T=S^\smile RR^\smile S=S^\smile S .
\]
Thus targets agree. For composable bridges $R:a\to b$,
$S:b\to c$, $(RS)(RS)^\smile
=R(SS^\smile)R^\smile=RK_bR^\smile=K_a$,
and exact convolution proves $RS$ is again a bridge.
The converse is a bridge in the opposite direction.
Only finitely many orbital supports occur, so the groupoid
is finite.

Every same-sort invariant relation equals its converse.
An isotropy arrow is therefore self-inverse, and its square
is $K_a$. A finite group of exponent two is elementary
abelian. If $R,S:a\to b$, then $RS^\smile$ is an isotropy
arrow sending $S$ to $R$, proving the torsor claim.
Finally an isotropy equivariant permutation of transitive
$M_a/K_a$ fixing a single quotient point fixes all points,
so must equal $K_a$. Thus different arrows cannot agree
at one point.
\end{proof}

\begin{lemma}[Transport of a dominant full exact involution]
\label{lem:foundational-rigid-gate-transport}
Let $\mathcal H$ be a rigid bridge groupoid as in
\cref{thm:foundational-rigid-groupoid}, with distinguished object
$(a_*,K_*)$ and a nonidentity isotropy arrow
$t_*\in\mathcal H(a_*,a_*)$. Assume that every \emph{original}
full exact invariant equivalence $Z_*>K_*$ on sort $a_*$
contains $t_*$. For every reached object $(b,K_b)$ and every
bridge $R:a_*\to b$ define
\[
                         t_b=R^\smile;t_*;R.
\]
Conjugation $Z_b\mapsto R;Z_b;R^\smile$ is an order
isomorphism from the poset of original full exact invariant
equivalences containing $K_b$ to the corresponding poset
above $K_*$, with inverse $Z_*\mapsto R^\smile;Z_*;R$.
Moreover $t_b$ is independent of $R$, is a nonidentity full exact
isotropy arrow with $t_b;t_b=K_b$, and is contained in every
\emph{original} full exact invariant equivalence $Z_b>K_b$.
Consequently its quotient graph $\bar t_b$ is contained in every
nonidentity \emph{full exact} equivalence on $M_b/K_b$.
No assumption of SAP or descent of arbitrary old exact domains
into quotient word-exact domains is needed.
\end{lemma}
\begin{proof}
If $R,S:a_*\to b$ are bridges, put $h=S;R^\smile$, an
isotropy arrow at $a_*$. Then $S=h;R$ and, since that isotropy
group is elementary abelian $2$,
\[
 S^\smile;t_*;S=R^\smile;h^\smile;t_*;h;R
                  =R^\smile;t_*;R.
\]
Composition and converse in the full exact bridge groupoid
show that the latter is a full exact isotropy arrow at $b$;
it is nonidentity because conjugation by the induced quotient
bijection is injective, and its square is $K_b$.

Now let $Z_b>K_b$ be any original full exact equivalence
on sort $b$. The old relation
\[
                     Z_*=R;Z_b;R^\smile
\]
is nonempty full exact by represented convolution. It is a
symmetric equivalence containing $K_*$: using
$R^\smile R=K_b$, $RR^\smile=K_*$ and
$K_b Z_b=Z_b=Z_b^2$, its square is itself.
Moreover $Z_*>K_*$, because $R$ induces a bijection
$M_{a_*}/K_*\to M_b/K_b$. More precisely, composition
with $R^\smile,R$ sends every full exact equivalence above
$K_*$ back to one above $K_b$ and is inverse to the first
conjugation, since
\[
 R^\smile;(R;Z_b;R^\smile);R=K_b;Z_b;K_b=Z_b
\]
and likewise in the other direction. Both maps preserve
inclusions, proving the asserted order isomorphism.
By the assumed dominance of $t_*$,
\[
 t_b=R^\smile;t_*;R\subseteq
 R^\smile;Z_*;R=K_b;Z_b;K_b=Z_b.
\]
Finally, for a full exact quotient equivalence
$\bar Z_b>\Id_{b/K_b}$, the inverse image $Z_b$
is an \emph{original} full exact equivalence strictly above
$K_b$ by the full exact quotient-word pullback in
\cref{lem:typed-simultaneous-blocks}. Applying the old-level
claim and projecting gives $\bar t_b\subseteq\bar Z_b$.
No reverse exact descent is invoked.
\end{proof}

\begin{lemma}[Exact bridges descend and lift across exact kernels]
\label{lem:foundational-rigid-quotient-bridge}
Let $E_a\le K_a$ and $E_b\le K_b$ be
full exact invariant equivalences on two
original named transitive sorts. Form their
actual simultaneous full exact quotient
$Q$ by $E_a,E_b$ (and any specified exact
kernels on the remaining original sorts).
Assume moreover that the projected
relations
\[
 \bar K_a=K_a/E_a,\qquad \bar K_b=K_b/E_b
\]
are \emph{full exact equivalences in the
quotient represented algebra}.
Then the map $R\mapsto\bar R$ taking
$E_a,E_b$-bisaturated original relations
to their quotient relations is a bijection
between the exact functional bridges
$R:(a,K_a)\to(b,K_b)$ and the exact
functional bridges
$\bar R:(Q_a,\bar K_a)\to(Q_b,\bar K_b)$.
It preserves converse and composition.
In particular no extra rigid bridge appears
in the quotient, and none disappears.

The full-exactness of the \emph{projected
gates} $\bar K_a,\bar K_b$ is an explicit
hypothesis: the general typed quotient
word-pullback theorem alone only proves
exactness in the reverse direction.
\end{lemma}
\begin{proof}
A $K_a,K_b$-saturated functional bridge
$R$ is automatically $E_a,E_b$-bisaturated,
and its quotient preserves the bridge
equations by the actual quotient product
formula
\cref{lem:typed-simultaneous-blocks}.
To verify \emph{quotient exactness},
fix any old atomic colour $c\subseteq R$.
The old equivariant quotient bijection
$M_a/K_a\to M_b/K_b$ is transitive on
its graph; hence all corresponding old
pairs lie in
\[
                   R=K_a;\{c\};K_b.
\]
Indeed the orbital $c$ meets every
corresponding pair of quotient $K$-classes,
and saturation by $K_a,K_b$ fills the whole
rectangular block. Passing to $Q$ gives
\[
             \bar R=\bar K_a;\bar c;\bar K_b,
\]
where $\bar c=E_a;\{c\};E_b$ is one
named actual quotient orbital. Since
$\bar K_a,\bar K_b$ and $\bar c$ are
full exact, quotient exact convolution
\cref{thm:full-typed-exact-balls}(ii)
makes $\bar R$ full exact.

Conversely, if $\bar R$ is a quotient
full exact bridge, its old inverse image
$R$ is original full exact by the
quotient full-word pullback of
\cref{lem:typed-simultaneous-blocks}.
The quotient bridge equalities pull
back to $RR^\smile=K_a$ and
$R^\smile R=K_b$ using the saturated
quotient product identity; $R$ is
thus an original exact functional bridge.
Inverse image and quotient are inverse
on all endpoint-saturated relations.
They commute with converse and with
composition by the same quotient product
formula. This establishes the bijection.
\end{proof}

\begin{lemma}[Affine exact rigid slices]
\label{lem:foundational-rigid-affine}
Let $D:a\to b$ be nonempty full exact and
$K_a,K_b$-bisaturated, supported entirely on the union
of the rigid arrows in a nonempty hom-set above.
Then $D$ is a union of \emph{whole} rigid arrows; for
some rigid arrow $d:a\to b$ and an exact subgroup
$S\le\mathcal H(b,b)$ it has the form $D=dS$.
Here $dS$ means the union of represented graphs of
the composite arrows, not an individual original atom.
\end{lemma}
\begin{proof}
The rigid arrow graph is one actual quotient orbital:
the quotient source is transitive and the arrow equivariant.
Distinct arrows are disjoint by freeness. Thus $D$,
being invariant and bisaturated, is a union of whole arrows.
Identify its nonempty support with a subset $A$ of the
elementary abelian isotropy group using one arrow $d$.
Translation by isotropy arrows preserves full exactness
by convolution. Pairwise intersection of finitely many
translates of $A$ becomes pairwise intersection of full
exact relations; their Helly property
\cref{thm:full-typed-exact-balls} supplies a common arrow.
The affine-translate lemma
\cref{lem:foundational-affine-translate-helly}
makes $A$ a coset. Absorb its shift into $d$ so that
$A=S$ is a subgroup. Since the represented product
$D^\smile D$ is the union of the arrows in $S$,
it is full exact; thus $S$ is an exact subgroup.
\end{proof}

\begin{lemma}[Rigid-product factorisation]
\label{lem:foundational-rigid-factor}
Suppose nonempty full exact $A:a\to b$ and $B:b\to c$
satisfy $AB\subseteq R$ for one full exact bridge
$R:a\to c$ from $(a,K_a)$. Then there are bridges
$R_1:a\to b$, $R_2:b\to c$ with
\[
 A\subseteq R_1,\quad B\subseteq R_2,\quad R_1R_2=R.
\]
This uses represented seriality, not SAP.
\end{lemma}
\begin{proof}
Every nonempty invariant relation on transitive source and
target sorts is serial and surjective. In particular
$\Id_b\subseteq BB^\smile$, whence
\[
 AA^\smile\subseteq ABB^\smile A^\smile
                 \subseteq RR^\smile=K_a.
\]
Set $R_1=K_aA$. It is full exact and left saturated.
The preceding inclusion gives $R_1R_1^\smile\subseteq K_a$;
seriality of $A$ gives the reverse inclusion.
Thus $R_1$ is a bridge; let $K_b=R_1^\smile R_1$.
Because $R_1B\subseteq K_aR=R$, the full exact relation
$B'=K_bB$ lies in $C=R_1^\smile R$, which is a
bridge from $(b,K_b)$ to $c$.
For any $(y,z)\in C$, surjectivity of $B'$ supplies
$y'$ with $(y',z)\in B'\subseteq C$.
Since $CC^\smile=K_b$, we have $yK_by'$,
and $K_bB'=B'$ gives $(y,z)\in B'$.
Hence $B'=C$. Put $R_2=C$.
Finally $R_1R_2=R_1R_1^\smile R=K_aR=R$.
\end{proof}

\begin{theorem}[Binary groupoid characters and phase marking]
\label{thm:foundational-rigid-character}
Assume the groupoid of
\cref{thm:foundational-rigid-groupoid}
has a distinguished nonidentity isotropy arrow
$t_*\in\mathcal H(a_*,a_*)$.
There exists a functor
$\chi:\mathcal H\to\mathbb F_2$
with $\chi(t_*)=1$. Every object $b$ has a
path-independent transported involution
$t_b=R^\smile t_*R\ne K_b$, with $\chi(t_b)=1$.
There is a global bit function on all reached quotient
points satisfying
\[
       \lambda(Rq)=\lambda(q)+\chi(R).
\]
Rigid quotient pair orbitals have one forced phase
difference; every nonrigid quotient pair orbital admits
both phase differences by appropriate global choices.
Thus the finite unary phase marking is reduct-surjective.
This alone does \emph{not} imply AP, SAP or CCO of the
marked age.
\end{theorem}
\begin{proof}
Choose a linear character of the finite source isotropy
$\mathbb F_2$-vector space taking $t_*$ to one.
Choose bridges $B_b:a_*\to b$, with $B_{a_*}=K_*$.
Every arrow $R:b\to c$ has a unique expression
$B_b^\smile hB_c$ with source isotropy $h$.
Set $\chi(R)=\chi_*(h)$; the torsor law makes it
functorial. Any two paths to one object differ by
isotropy conjugation, and that group is abelian,
so transporting $t_*$ is path-independent, nontrivial,
and has character one.

The groupoid acts freely on the disjoint union of its
quotient point sets. Choose one initial bit for each
orbit and propagate by the displayed character equation.
Freeness makes this well-defined. Each rigid arrow is
one equivariant quotient orbital and fixes its pairwise
phase difference. If a pair orbital is not rigid, its
two endpoints belong to different groupoid orbits:
otherwise some rigid arrow would carry one point to
the other, and their actual quotient orbital would
be exactly that rigid arrow. Initial bits on different
orbits may be toggled independently, realizing both
differences. Pulling back the global bits to all
original real points gives the asserted markings.
\end{proof}

\begin{lemma}[Finite phase constraints extend by groupoid seeds]
\label{lem:foundational-finite-phase-extension}
Use a free action on quotient points of a finite connected
rigid groupoid with a fixed character
$\chi:\mathcal H\to\mathbb F_2$ as above. Let $F$ be any
finite set of quotient points and let
$\epsilon:F\to\mathbb F_2$ be prescribed bits.
Then there is a \emph{global} phase marking $\lambda$
satisfying
\[
 \lambda(Rq)=\lambda(q)+\chi(R)
 \quad\text{for all rigid arrows $R$},
 \qquad \lambda|_F=\epsilon ,
\]
if and only if, whenever $q,r\in F$ belong to the
same rigid-groupoid orbit and $Rq=r$, one has
\[
             \epsilon(r)=\epsilon(q)+\chi(R).
 \tag{FS}\label{eq:foundational-phase-seed-test}
\]
It suffices to check the prescribed bits pairwise; no
additional higher-arity character holonomy condition is
needed. This does \emph{not} imply a finite projection
theorem for moving-owner pointwise kernels from
\cref{thm:finite-fibre-kernel-csp}.
\end{lemma}
\begin{proof}
Necessity follows from the defining character equation.
For sufficiency, split the entire quotient-point universe
into its groupoid orbits. On each orbit meeting $F$ choose
one $q_0\in F$ and prescribe $\lambda(q_0)=\epsilon(q_0)$;
on all other orbits choose an arbitrary seed.
Define $\lambda(Rq_0)=\epsilon(q_0)+\chi(R)$.
Freeness in
\cref{thm:foundational-rigid-groupoid}
makes the transporting arrow unique, and
functoriality of $\chi$ proves every required
equation. For any other $r\in F$ in that orbit,
the pairwise condition \eqref{eq:foundational-phase-seed-test}
gives the desired value $\epsilon(r)$.
There are no constraints between different orbits.
\end{proof}
